\section{Related Work}
\label{sec:related}

A wave of recent benchmarks evaluates LLM agents on security tasks, summarized in \cref{tab:related}. \textsc{SWE-bench}~\cite{jimenez2024swebench} tests bug fixing on Python repositories; \textsc{CVE-Bench}~\cite{cvebench} targets web-application exploitation in containerized apps; \textsc{CyberGym}~\cite{wang2025cybergym} and \textsc{CyberGym-E2E}~\cite{cybergyme2e} evaluate PoC reproduction and end-to-end discovery/patching on C/C++ projects sourced from OSS-Fuzz; \textsc{BountyBench}~\cite{zhang2025bountybench} spans the offense--defense lifecycle with real bug-bounty dollar values; \textsc{SEC-bench}~\cite{secbench} automates benchmark construction from CVE records; \textsc{ExploitGym}~\cite{exploitgym} and \textsc{ExploitBench}~\cite{exploitbench} grade exploitation across userspace, V8, and kernel targets; and \textsc{K-Repro}~\cite{krepror} studies Linux-kernel N-day reproduction from patches. What unites these benchmarks is what they hand the agent: source code, a build script, or a pre-built executable. None requires the agent to acquire firmware, extract a filesystem from a proprietary container, or rehost a cross-architecture binary---the steps that dominate real IoT reproduction and where, as we show, agents fail catastrophically.

\begin{table*}[t]
    \centering
    \caption{Positioning of \sys against representative agentic security benchmarks. \emph{Artifact} is the starting material handed to the agent; \emph{X-arch} indicates whether non-x86 targets are involved; \emph{Rehost} indicates whether the agent must rehost/emulate the target to reproduce; \emph{Graded} indicates partial-credit (non-binary) scoring.}
    \label{tab:related}
    \small
    \begin{tabular}{lllcccc}
        \toprule
        \textbf{Benchmark} & \textbf{Domain} & \textbf{Starting artifact} & \textbf{X-arch} & \textbf{Rehost} & \textbf{Graded} & \textbf{\# Tasks} \\
        \midrule
        SWE-bench~\cite{jimenez2024swebench}      & Python repos      & Source + issue           & \cross & \cross & \cross & 2294 \\
        CVE-Bench~\cite{cvebench}                 & Web apps          & Running container        & \cross & \cross & \cross & 40 \\
        CyberGym~\cite{wang2025cybergym}          & C/C++ (OSS-Fuzz)  & Source + harness         & \cross & \cross & \cross & 1507 \\
        CyberGym-E2E~\cite{cybergyme2e}           & C/C++ (OSS-Fuzz)  & Source + build env       & \cross & \cross & \cross & 920 \\
        BountyBench~\cite{zhang2025bountybench}   & Web/OSS systems   & Running system           & \cross & \cross & \cross & 40 \\
        SEC-bench~\cite{secbench}                 & C/C++             & Dockerized source        & \cross & \cross & \cross & 200 \\
        ExploitGym~\cite{exploitgym}              & Userspace/V8/krnl & Binary + PoV             & \cross & \cross & \tick  & 898 \\
        ExploitBench~\cite{exploitbench}          & V8 engine         & Built d8 binary          & \cross & \cross & \tick  & 41 \\
        K-Repro~\cite{krepror}                    & Linux kernel      & Patch + buildable kernel & \cross & \tick  & \cross & 100 \\
        \midrule
        \textbf{\sys (ours)}                      & \textbf{IoT firmware} & \textbf{CVE number only} & \tick & \tick & \tick & \textbf{124} \\
        \bottomrule
    \end{tabular}
\end{table*}

The firmware-analysis community has spent over a decade automating exactly the steps that agents stumble on. \textsc{Firmadyne}~\cite{firmadyne} pioneered full-system emulation with kernel replacement, achieving IP connectivity on about 21\% of firmware images; \textsc{FirmAE}~\cite{firmae} raised this to 79\% through arbitration techniques that detect and work around boot failures; \textsc{Greenhouse}~\cite{greenhouse} introduced user-space single-service rehosting, handling 40\% of firmware without booting a full system; \textsc{FIRMWELL}~\cite{firmwell} achieved 68\% with dependency-aware multi-process emulation; \textsc{Penguin}~\cite{penguin} uses target-centric configuration files to automate NVRAM and peripheral setup; and \textsc{Pandawan}~\cite{pandawan} augments kernels for module analysis. These tools are mature, openly available, and installable via standard package managers---yet, as our experiments show, LLM agents given \texttt{sudo} access and a 30-minute budget almost never use them effectively. That gap between tool availability and tool orchestration is precisely what \sys is designed to measure.

A separate strand of work evaluates graded rather than binary outcomes. \textsc{ExploitBench}~\cite{exploitbench} decomposes exploitation into a 16-flag capability ladder from crash to arbitrary code execution, and \textsc{CyberGym-E2E}~\cite{cybergyme2e} uses staged validation (S1--S4) that credits partial progress through discovery, patching, and functionality testing. Both demonstrate the diagnostic value of partial credit, but neither targets a pipeline where the dominant difficulty lies \emph{before} the vulnerability is even reachable. The \rrs extends this philosophy to the IoT pipeline, where acquisition, extraction, and rehosting gate all downstream tasks, and makes the dependency structure explicit so that an agent cannot earn exploitation credit by crashing a program it wrote itself. We also note detection-oriented benchmarks such as \textsc{SecLLMHolmes}~\cite{secllmholmes} and \textsc{AutoAdvExBench}~\cite{autoadvexbench}, which assess reasoning about vulnerabilities rather than reproduction of them; \sys is complementary, measuring whether an agent can turn reasoning into a verified trigger against real firmware.
